% ============================================================================
% chapters/ch5_evaluation.tex
% ============================================================================
\chapter{Evaluation}
\label{ch:evaluation}

% ============================================================================
\section{Experimental Setup}
\label{sec:experimental_setup}
% ============================================================================

All experiments are run on the 26 Garmin Decoded recording sessions described
in \Cref{sec:dataset}.
Experiments are identified by artefact run codes (A1--A5) that correspond to
MLflow run IDs stored in \texttt{mlruns/} and artefact directories under
\texttt{artifacts/}.
All results are reproducible via \texttt{python run\_pipeline.py} with the
command recorded in each artefact's \texttt{run\_config.json}.

\begin{table}[ht]
  \centering
  \caption{Ablation experiment definitions.}
  \label{tab:ablation_experiments}
  \begin{tabularx}{\linewidth}{clX}
    \toprule
    \textbf{ID} & \textbf{Name} & \textbf{Description} \\
    \midrule
    A1 & Baseline inference
       & Pre-trained model, no adaptation, raw production data. \\
    A2 & Baseline (converted)
       & Pre-trained model, no adaptation, unit-converted data. \\
    A3 & Fine-tuned
       & Model fine-tuned on labelled wrist data; best fold from 5-fold CV. \\
    A4 & AdaBN+TENT
       & AdaBN followed by TENT on unlabelled production data. \\
    A5 & Pseudo-label
       & Curriculum pseudo-labelling with EWC, 9 adaptation stages. \\
    \bottomrule
  \end{tabularx}
\end{table}

\section{Evaluation Metrics}

\begin{description}
  \item[Classification accuracy] fraction of windows with correct top-1
    prediction on a held-out source-domain test set.
  \item[Macro F1] unweighted mean of per-class F1 scores; accounts for class
    imbalance.
  \item[Mean confidence $\bar{c}$] mean maximum softmax probability across all
    production windows; primary monitoring signal.
  \item[ECE] Expected Calibration Error (Equation~\ref{eq:ece}); measures
    calibration quality before and after temperature scaling.
  \item[Transition rate] fraction of consecutive window pairs with different
    predicted class; measures temporal coherence.
\end{description}

% ============================================================================
\section{Classification Performance}
\label{sec:classification_performance}
% ============================================================================

\begin{table}[ht]
  \centering
  \caption{Classification results across ablation experiments.
           \textsuperscript{†}Production data without unit conversion.}
  \label{tab:results}
  \begin{tabular}{clccc}
    \toprule
    \textbf{ID} & \textbf{Experiment}
      & \textbf{Val Acc (\%)} & \textbf{Macro F1}
      & \textbf{$\bar{c}$} \\
    \midrule
    A1 & Baseline (raw)\textsuperscript{†} & 14.52 & \todo{TBD} & \todo{TBD} \\
    A2 & Baseline (converted)              & \todo{TBD} & \todo{TBD} & \todo{TBD} \\
    A3 & Fine-tuned                        & 89.90 & \todo{TBD} & 0.899 \\
    A4 & AdaBN+TENT                        & \todo{TBD} & \todo{TBD} & 0.823 \\
    A5 & Pseudo-label                      & 96.90 & \todo{TBD} & \todo{TBD} \\
    \bottomrule
  \end{tabular}
\end{table}

\paragraph{A1 vs A2: Unit conversion impact.}
The accelerometer unit mismatch caused cross-user accuracy to fall below the
random-chance baseline of 9.09\,\% for 11 classes.
After applying the milli-g to m/s\textsuperscript{2} conversion, accuracy
recovered to the levels reported for A2.
This experiment demonstrates that unit-compatibility validation is a
prerequisite for any meaningful HAR evaluation.

\paragraph{A3: Fine-tuned baseline.}
The fine-tuned model achieves 89.9\,\% validation accuracy with mean confidence
$\bar{c} = 0.899$, establishing the performance target for the adaptation
strategies.

\paragraph{A4: AdaBN+TENT.}
The adaptation run on 2026-02-23 resulted in mean confidence
$\bar{c} = 0.823$ ($\Delta\bar{c} = -0.076$), exceeding the 5\,\% canary
tolerance ($-0.045$).
The adapted model was archived; the production model remains A3.
The TENT rollback gate was not triggered (entropy did not increase), indicating
that the confidence drop was genuine and not an artefact of adaptation instability.

\paragraph{A5: Pseudo-labelling.}
The pseudo-label run (artefact \texttt{20260219\_134233}) completes 9
adaptation stages, achieves in-loop validation accuracy 96.9\,\%, and passes
12 of 12 audit checks.
Full confusion matrix and per-class F1 scores are provided in
\Cref{app:full_results}.

% ============================================================================
\section{Calibration Results}
\label{sec:calibration_results}
% ============================================================================

Temperature scaling was applied to all experiment variants.
\todo{Insert ECE before/after table after running calibration stage on all
variants.}

% ============================================================================
\section{Drift Detection Accuracy}
\label{sec:drift_detection}
% ============================================================================

\todo{Fill in after running batch drift analysis across 26 sessions via
\texttt{python scripts/analyze\_drift\_across\_datasets.py}.}

The drift detection framework is validated by replaying the A1 scenario
(raw production data) through the monitoring pipeline and verifying that:
\begin{itemize}
  \item Layer~3 z-score alert fires for the $a_z$ channel
        ($z_{a_z} \approx (−1001.6 − (−3.5))/3.2 \approx 312$).
  \item PSI alert fires with aggregated PSI $\gg 1.50$ (critical threshold).
  \item The trigger engine votes CRITICAL and initiates a retraining event.
\end{itemize}

% ============================================================================
\section{Robustness Testing}
\label{sec:robustness}
% ============================================================================

\todo{Run robustness tests via \texttt{src/robustness.py} after batch run
completes.}

The robustness module (\texttt{src/robustness.py}, 449 lines) injects five
noise types --- Gaussian noise, bias shift, axis dropout, random spike
injection, and temporal jitter --- at three severity levels (low, medium, high).
Results will show how monitoring Layer~1 confidence responds to each noise type
and at what severity level the trigger fires.
